\section{Introduction}\label{sec:introduction}
An aircraft-inspection system needs more than a low-error pose. An observation must receive a fresh airframe-relative correction before its deadline; a continuing route needs acceptable error over a longer interval. The same localization computation can satisfy one requirement and fail the other. Moreover, a reliability statement established for one registration policy, sensor mixture or installation service does not automatically apply after a software update or a revision of the pose. The engineering decision is therefore which computation to execute, when to expose its result, and which evidence remains applicable to the result that a task actually consumes.

This problem connects engineering decision knowledge to an executable numerical service. Work on process and decision knowledge representation shows how explicit dependencies support engineering decisions \cite{wang2021process,hao2021knowledge,su2025manufacturing}. Contract-based design makes assumptions and guarantees part of a component interface \cite{benveniste2018contracts}. Anytime autonomous systems already connect computational contracts to estimation and control \cite{pant2021}; perception scheduling also studies the tradeoff between latency and precision \cite{aldana2023}. The unresolved issue addressed here is narrower: an interrupted registration produces correlated revisions of one source, whereas a population-risk statement refers to a complete policy, installation model and specified task event. Their identities and lifetimes must remain consistent from pre-sensing selection to installed pose and final evidence attachment.

Our previous A19/A20 development established retained registration prefixes, source-time replacement and empirical computation--error tradeoffs. It also exposed three limitations. A guarded CART scheduler often behaved like its flat predictor; fixed two-iteration registration could outperform it on ordinary complete routes; and metadata prototypes did not establish an integrated evidence lifecycle. Increasing a favorable mean or adding a knowledge-labelled diagram would not resolve these limitations. We retain the estimator and assets and instead test the contract mechanism against a task-aware, risk-calibrated procedural implementation supplied with the same information.

We represent sources, versions, policies, operating profiles, processing services, requests and calibrated event cards as typed entities with explicit dependencies and execution predicates. A task declared before sensing selects an admissible frozen policy or receives a refusal. During execution, budget and predicted-time constraints restrict computation, correlated results replace the same source posterior, and measured motion is replayed to installation. A source-scoped attachment monitor checks the identity and history needed to associate the installed result with its population-risk statement. Changes to a dependency revoke affected claims. A numerical result can remain available with a provisional status after its statistical attachment becomes inapplicable.

The study makes four connected contributions.
\begin{enumerate}
\item An executable representation of the relationship between engineering requirements, numerical versions and statistical evidence, including pre-sensing policy commitment, selective dependency invalidation and permanent source-level attachment refusal after a detected invalid version attempt. These mechanisms are evaluated separately from pose accuracy.
\item A common localization service that combines retained ICP computation, constrained publication, source replacement and causal replay. A dimensionally consistent work objective and two explicit timing ledgers permit comparisons with fixed, learned and earlier RAVC policies under the same resources.
\item A prospectively frozen empirical assessment of contract modules, task-driven selection, directly calibrated union events, certificate margin, unfitted geometry, binding time/profile constraints, continuous routes and influential failures. A fully validated procedural comparator prevents the engineering representation from receiving credit for information or checks withheld from its baseline.
\item An extension of the same operating-contract interface to image acquisition, capture-time association and report delivery. Requirement containment and event-specific producer dependencies determine evidence reuse or renewal. Two further executed studies compare this extension with complete procedural semantics and cached direct recalibration, retaining their equal or superior results.
\end{enumerate}

The central hypothesis is that executable knowledge can prevent an engineering task from using evidence outside its scope while making computational tradeoffs explicit. It does not require graph storage to outperform a semantically equivalent program in localization accuracy. We accordingly separate numerical localization, engineering decision quality and evidence-attachment correctness. The localization and task-extension experiments use simulated geometric sensors and a declared processing model; original platform photographs and previous studies remain as documented context. No new physical localization experiment is claimed.

Figure~\ref{fig:contracts} follows the complete decision path. Sections~\ref{sec:problem}--\ref{sec:method} define its objects, rules and execution. Section~\ref{sec:theory} states the properties needed to justify replacement and evidence applicability. Sections~\ref{sec:design} and \ref{sec:evaluation} report the localization and route study. Sections~\ref{sec:r9-executor}--\ref{sec:r9-theory} extend its task interface; Sections~\ref{sec:r9-design} and \ref{sec:r9-results} evaluate that extension on separately generated inputs. Each evidence block retains its own sampling units, calibration family and timing boundary.
\begin{figure}[t]\centering
\includegraphics[width=\linewidth]{r8/architecture.pdf}
\caption{Executable engineering requirements, numerical execution and evidence applicability. (a) An exogenous event requirement, risk allowance, deadline and budget constrain a dependency-bound policy catalogue. (b) Retained computation revises one immutable source; ordered replacement and measured-motion replay produce the live pose. The airframe surface and source points illustrate an actual excluded pilot input. (c) Final attachment checks join an event card to the active source version and its service; invalid history or dependencies remove the bound. The feedback arrow denotes the interface for a future request or recalibration; autonomous feedback replanning and a physical control experiment are not evaluated. OpenAI ChatGPT/Codex (GPT-6) assisted the archived Python schematic code.}\label{fig:contracts}
\end{figure}
